% !TeX root = ../../Book4.tex
% 纲要标题：### 15.3 Grothendieck 谱序列
% 纲要位置：计划纲要.md，第 2875 行。

\section{Grothendieck 谱序列}
\label{b4:sec:1503}

复合两个左正合函子时，第一个函子的导出群不能一般地丢掉。内射像对第二个函子的无上同调条件，使这组中间信息组织成收敛到复合导出的谱序列。

% 纲要内容：
%
% * 左正合函子复合
% * 内射映到后一函子无上同调对象的条件
% * 谱序列及低阶正合列
%

\subsection{左正合函子的复合}
\label{b4:subsec:1503:01}

令 \(F:\mathcal A\to\mathcal B\)、\(G:\mathcal B\to\mathcal C\) 为加性左正合函子，前两个交换范畴有足够内射对象。对 \(M\in\mathcal A\) 取内射消解 \(M\to I^\bullet\)，\(F(I^\bullet)\) 的上同调是 \(R^qF(M)\)。因此它通常不是 \(F(M)\) 的正合消解，而是一个必须保留正次上同调的复形。上一节正好提供了对这种复形再作用 \(G\) 的方法。

\ProseParagraphBreak
第二套超上同调谱序列的第二页已经是 \(R^pG(R^qF(M))\)。尚待解决的只有目标：\(\mathbb H^*(G,F(I^\bullet))\) 是否就是复合函子的导出？这一识别需要各项 \(F(I^a)\) 对 \(G\) 无上同调，从而第一套谱序列退化。这解释了复合函子定理中看似技术性的对象条件。

\subsection{内射对象的无上同调像}
\label{b4:subsec:1503:02}

\begin{theorem}[Grothendieck 谱序列]\label{b4:thm:grothendieck-spectral-sequence}
设 \(\mathcal A,\mathcal B,\mathcal C\) 为交换范畴，\(\mathcal A,\mathcal B\) 有足够内射对象。设 \(F:\mathcal A\to\mathcal B\)、\(G:\mathcal B\to\mathcal C\) 为加性左正合函子，并且
\[
 R^pG(F(I))=0\quad(p>0)
\]
对每个内射对象 \(I\in\mathcal A\) 成立。则对每个 \(M\in\mathcal A\)，有自然的第一象限强收敛\term[grothendieck-puxulie]{Grothendieck 谱序列}[Grothendieck spectral sequence]
\[
 E_2^{p,q}=R^pG(R^qF(M))\Longrightarrow R^{p+q}(GF)(M).
\]
\end{theorem}
\begin{proof}
取 \(M\to I^\bullet\) 的内射消解，令 \(C=F(I^\bullet)\)。它位于非负次数。取\cref{b4:thm:ce-resolution}构造的 \(C\) 的 Cartan--Eilenberg 消解 \(C\to J^{\bullet,\bullet}\)。对 \(G(J)\) 的两种滤过应用\cref{b4:thm:hypercohomology-spectral-sequences}。第一套给出
\[
 {}^{\mathrm I}E_2^{a,b}=H^a(R^bG(F(I^\bullet))).
\]
假设使 \(b>0\) 全部为零，而零行是 \(H^a(GF(I^\bullet))\)。因此边缘同态自然识别
\[
 H^n\operatorname{Tot}(G(J))
 \simeq H^n(GF(I^\bullet))=R^n(GF)(M).
\]
第二套则给出
\[
 {}^{\mathrm {II}}E_2^{p,q}=R^pG(H^q(F(I^\bullet)))
 =R^pG(R^qF(M)),
\]
目标是刚识别的同一总上同调。两个消解次数非负，所以滤过逐度有限，这证明强收敛。

对 \(M\to N\)，内射比较定理给出 \(I_M\to I_N\)，再用 Cartan--Eilenberg 比较定理得到双复形映射。第二页是导出函子的自然映射。两种内射提升链同伦，故在 \(H^qF(I)\) 上相同，因而第二页相同；以后各页随之相同。目标的识别来自增广 \(G(C)\to\operatorname{Tot}(G(J))\)，与这些比较映射相容，而 \(GF(I)\) 上的同伦给出通常导出函子的自然态射。于是谱序列及其收敛识别都与选择无关。
\end{proof}

若 \(F\) 是一个正合函子的右伴随，它保持内射对象，自动满足定理条件；这是\cref{b4:prop:right-adjoint-injectives}给出的便于检验的充分条件。保持内射不是必要条件，只须高次 \(R^pG(F(I))\) 消失。若 \(F\) 正合，谱序列只剩底行，恢复\cref{b4:thm:derived-composite-exact-first}；若 \(G\) 正合，只剩最左列，恢复\cref{b4:prop:derived-composite-exact-second}。

\subsection{谱序列与低阶正合列}
\label{b4:subsec:1503:03}

\begin{proposition}\label{b4:prop:spectral-five-term}
设交换范畴中的第一象限上同调谱序列从第二页起给出，并强收敛于带有限滤过的 \(H^*\)。则有自然正合列
\[
 0\longrightarrow E_2^{1,0}\longrightarrow H^1
 \longrightarrow E_2^{0,1}\xrightarrow{\mathrm{d}_2}E_2^{2,0}
 \longrightarrow H^2.
\]
最右端不宣称为满态射。
\end{proposition}
\begin{proof}
位置 \((1,0)\) 没有第 \(r\ge2\) 页的入箭头或出箭头，所以 \(E_\infty^{1,0}=E_2^{1,0}\)。位置 \((0,1)\) 只有一条可能的箭头 \(\mathrm{d}_2\) 指向 \((2,0)\)，所以 \(E_\infty^{0,1}=\ker \mathrm{d}_2\)。总次数一的滤过给出 \(0\to E_\infty^{1,0}\to H^1\to E_\infty^{0,1}\to0\)，再包含到 \(E_2^{0,1}\)，得到前半列的正合性。

位置 \((2,0)\) 没有出箭头，唯一可能的入箭头也是这条 \(\mathrm{d}_2\)，故 \(E_\infty^{2,0}=\operatorname{coker}\mathrm{d}_2\)。总次数二的最深非零层为 \(F^2H^2\)，并有 \(F^3H^2=0\)，所以 \(E_\infty^{2,0}=F^2H^2\hookrightarrow H^2\)。这给出余下正合性。全部态射来自滤过及谱序列的自然态射，因此自然。
\end{proof}

代入 Grothendieck 谱序列，得到
\[
\begin{aligned}
0&\longrightarrow R^1G(FM)\longrightarrow R^1(GF)(M)
\longrightarrow G(R^1F(M))\\
&\xrightarrow{\mathrm{d}_2}R^2G(FM)\longrightarrow R^2(GF)(M).
\end{aligned}
\]
这个列说明 \(R^1(GF)\) 同时接受后一函子的一次导出和前一函子的一次导出的贡献，而中间的 \(\mathrm{d}_2\) 检测两种计算能否相容。定理及低阶列可对照 Weibel\cite[Cohomological Setup 5.8.1, Grothendieck Spectral Sequence Theorem 5.8.3]{Weibel1994HomologicalAlgebra}；该书函子的字母次序与这里相反，阅读时应按复合次序识别。
